\subsection{Temporal Models}
\label{sec:temporal-models}
\inlinenotation Linearization typically written as $\delta X$.

Model \bi{robot dyn} as \bi{Cont-time non-lin. system of ODE}:\\
$\dot{\vec{x}} = \vec{f}_C(\vec{x}(t), \vec{u}(t), \vec{w}(t))$, measurement $\vec{z}(t) = \vec{h}(\vec{x}(t)) + \vec{v}(t)$.

With: $\pardiff{t}\vec{x}(t) = f_C(\vec{x}(t), \vec{u}(t))$ the model for the robot state update and $\vec{h}(\vec{x}(t))$ the model for the measurements (e.g. for IMU)

\vspace{0.5mm}
\hrule
\vspace{0.5mm}

\bi{Linearised} at $\vec{f}_C(\vec{\overline{x}}, \vec{\overline{y}}) = 0$, at \bi{equilibrium} {\scriptsize (if $\neq 0$, then add it)}:\\
$\delta \vec{\dot{x}}(t) = \mat{F}_C \delta \vec{x}(t) + \mat{G}_C \delta \vec{u}(t) + \mat{L}_C \vec{w}(t)$;
$\delta \vec{z}(t) = \mat{H} \delta \vec{x}(t) + \vec{v}(t)$.

$\mat{F}_C$ system (often a Jac. / Taylor, vars are $x_i$. If no $x$ in $\mat{F}$, then linear),
$\mat{G}$ input gain (often Jac / Taylor, vars are $u_i$, if no $u_i$, then lin.), $\vec{w}$ process noise, $\mat{H}$ is measurement, $\vec{v}$ measurement noise,
both zero-mean \bi{Gauss. White Noise Proc.}. $\vec{u}_k$ inputs at time $k$.

% TODO: Taylor approximation
\vspace{0.5mm}
\hrule
\vspace{0.5mm}


\bi{Discretized}
$\vec{x}_k = \vec{f}(\vec{x}_{k - 1}, \vec{u}_k, \vec{w}_k)$;
$\vec{z}_k = \vec{h}(\vec{x}_k) + \vec{v}_k$

\bi{Linearised}:
$\delta \vec{x}_k = \mat{F} \delta \vec{x}_{k - 1} + \mat{G}_k \delta \vec{u}_k + \mat{L}_k \vec{w}_k$;
$\delta \vec{z}_k = \mat{H}_k \delta \vec{x}_k$

For discretized, iterative integral from $t_{k - 1}$ to $t_k$

Linearization happens typically with one of the below:

\shortdefinition[Euler-Forward] $\vec{x}_k = \vec{x}_{k - 1} + \Delta t \vec{f}_C(\vec{x}_{k - 1}, \vec{u}_{k - 1}, t_{k - 1})$

\newpage
\shortdefinition[Trapezoidal num. int]
$\Delta \vec{x}_1 = \Delta t \vec{f}_C (\vec{x}_{k - 1}, \vec{u}_{k - 1}, t_{k - 1})$\\
$\Delta \vec{x}_2 = \Delta t \vec{f}_C (\vec{x}_{k - 1} + \Delta \vec{x}_1, \vec{u}_{k}, t_{k})$, then:\\
$\vec{x}_k = \vec{x}_{k - 1} + 0.5 \cdot (\Delta \vec{x}_1 + \Delta \vec{x}_2)$
